\documentclass[10pt,a4paper]{amsart}
\usepackage{amsmath,amssymb,amsthm,graphicx}
% Retained minimal preamble for the Stacks Project chapter sites-cohomology.tex (Cohomology on Sites).
% Verbatim excerpt of the chapter preamble: Xy-pic setup (native preamble lines 14-18) and the native
% theorem declarations with the mathematical macros (native preamble lines 154-217). Copied mechanically;
% nothing added, removed or redefined. The wrapper class and the amsmath/amssymb/amsthm/graphicx base
% packages are supplied by the assembler, so the class/externaldocument part of the native preamble is omitted.
% For commutative diagrams we use Xy-pic
\usepackage[all]{xy}

% We use 2cell for 2-commutative diagrams.
\xyoption{2cell}
% Theorem environments.
%
\theoremstyle{plain}
\newtheorem{theorem}[subsection]{Theorem}
\newtheorem{proposition}[subsection]{Proposition}
\newtheorem{lemma}[subsection]{Lemma}

\theoremstyle{definition}
\newtheorem{definition}[subsection]{Definition}
\newtheorem{example}[subsection]{Example}
\newtheorem{exercise}[subsection]{Exercise}
\newtheorem{situation}[subsection]{Situation}

\theoremstyle{remark}
\newtheorem{remark}[subsection]{Remark}
\newtheorem{remarks}[subsection]{Remarks}

\numberwithin{equation}{subsection}

% Macros
%
\def\lim{\mathop{\mathrm{lim}}\nolimits}
\def\colim{\mathop{\mathrm{colim}}\nolimits}
\def\Spec{\mathop{\mathrm{Spec}}}
\def\Hom{\mathop{\mathrm{Hom}}\nolimits}
\def\Ext{\mathop{\mathrm{Ext}}\nolimits}
\def\SheafHom{\mathop{\mathcal{H}\!\mathit{om}}\nolimits}
\def\SheafExt{\mathop{\mathcal{E}\!\mathit{xt}}\nolimits}
\def\Sch{\mathit{Sch}}
\def\Mor{\mathop{\mathrm{Mor}}\nolimits}
\def\Ob{\mathop{\mathrm{Ob}}\nolimits}
\def\Sh{\mathop{\mathit{Sh}}\nolimits}
\def\NL{\mathop{N\!L}\nolimits}
\def\CH{\mathop{\mathrm{CH}}\nolimits}
\def\proetale{{pro\text{-}\acute{e}tale}}
\def\etale{{\acute{e}tale}}
\def\QCoh{\mathit{QCoh}}
\def\Ker{\mathop{\mathrm{Ker}}}
\def\Im{\mathop{\mathrm{Im}}}
\def\Coker{\mathop{\mathrm{Coker}}}
\def\Coim{\mathop{\mathrm{Coim}}}

% Boxtimes
%
\DeclareMathSymbol{\boxtimes}{\mathbin}{AMSa}{"02}

%
% Macros for moduli stacks/spaces
%
\def\QCohstack{\mathcal{QC}\!\mathit{oh}}
\def\Cohstack{\mathcal{C}\!\mathit{oh}}
\def\Spacesstack{\mathcal{S}\!\mathit{paces}}
\def\Quotfunctor{\mathrm{Quot}}
\def\Hilbfunctor{\mathrm{Hilb}}
\def\Curvesstack{\mathcal{C}\!\mathit{urves}}
\def\Polarizedstack{\mathcal{P}\!\mathit{olarized}}
\def\Complexesstack{\mathcal{C}\!\mathit{omplexes}}
% \Pic is the operator that assigns to X its picard group, usage \Pic(X)
% \Picardstack_{X/B} denotes the Picard stack of X over B
% \Picardfunctor_{X/B} denotes the Picard functor of X over B
\def\Pic{\mathop{\mathrm{Pic}}\nolimits}
\def\Picardstack{\mathcal{P}\!\mathit{ic}}
\def\Picardfunctor{\mathrm{Pic}}
\def\Deformationcategory{\mathcal{D}\!\mathit{ef}}
\usepackage{cleveref}
\crefname{originalreference}{original reference}{original references}
\Crefname{originalreference}{Original reference}{Original references}
\makeatletter
\newlabel{lemma-perfect-precise}{{08CP}{source}{}{sourcepointer-1}{}}
\newlabel{lemma-perfect-precise@cref}{{[originalreference][0][]08CP}{[1][0][]source}{}{}{}}
\makeatother
\crefname{originalreference}{original reference}{original references}
\Crefname{originalreference}{Original reference}{Original references}
\makeatletter
\newlabel{lemma-pseudo-coherent-independent-representative}{{08CC}{source}{}{sourcepointer-2}{}}
\newlabel{lemma-pseudo-coherent-independent-representative@cref}{{[originalreference][0][]08CC}{[1][0][]source}{}{}{}}
\makeatother
\crefname{originalreference}{original reference}{original references}
\Crefname{originalreference}{Original reference}{Original references}
\makeatletter
\newlabel{lemma-tor-amplitude}{{0654}{source}{}{sourcepointer-3}{}}
\newlabel{lemma-tor-amplitude@cref}{{[originalreference][0][]0654}{[1][0][]source}{}{}{}}
\makeatother
\title{Perfect objects are exactly the pseudo-coherent objects of locally finite tor dimension}
\author{The Stacks Project Authors}
\date{stacks-project:sites-cohomology.tex:local-snapshot-20261011; GNU FDL 1.2 or later}
\begin{document}
\maketitle
\noindent\textit{Editorial note.} All delivered mathematical text is the author's own, copied from the pinned Stacks Project chapter file sites-cohomology.tex; no statement or proof step was rewritten, reordered or removed.\par\medskip
\noindent\textit{Editorial note.} The counted proof is the author's complete argument: the proof is written out in full: the local strictly perfect model for the first implication with pseudo-coherence by locality and flatness of direct summands, and the local tor amplitude bound with the precise criterion for the second.\par\medskip
\noindent\textit{Editorial note.} The label lemma-perfect names the counted statement, and the proof environment immediately following it in the source is delivered verbatim as its proof.\par\medskip
\noindent\textit{Editorial note.} The counted statement is carried by the author's own lemma environment in the chapter source.\par\medskip
\noindent\textit{Editorial note.} The statement and proof are copied from the pinned chapter file; the three cited inputs are kept as external pointers with their canonical Stacks tags.\par\medskip
\noindent\textit{Allowed standard prerequisites (editorial).}\begin{itemize}
\item Strictly perfect complexes as local models. Same chapter, the definition of a perfect object through local representatives by strictly perfect complexes. Supplies the starting point of the implication (1) implies (2).
\item Pseudo-coherence of strictly perfect representatives and flatness of direct summands. Same chapter, the locality lemma for pseudo-coherence and the tor amplitude lemma for direct summands of finite free modules. Supplies pseudo-coherence and finite tor dimension of the local models in the first implication.
\item The precise criterion for perfection. Same chapter, the lemma stating that tor amplitude in a bounded interval plus (a-1)-pseudo-coherence forces perfection. Supplies the second implication once the local tor amplitude bounds are available.
\end{itemize}\medskip
\ifdefined\hypertarget\hypertarget{sourcepointer-1}{}\fi
\noindent\textit{Source reference.} Original reference 08CP: sources/stacks-sites-cohomology/pinned-original-source.tex, source lines 12774--12774. Prerequisite pointer invoked at proof line(s) 12843; it names a standard fact of the same chapter that the counted argument uses, and every step of the counted proof built on it is retained verbatim in the delivered body.\par\medskip
\ifdefined\hypertarget\hypertarget{sourcepointer-2}{}\fi
\noindent\textit{Source reference.} Original reference 08CC: sources/stacks-sites-cohomology/pinned-original-source.tex, source lines 12066--12066. Prerequisite pointer invoked at proof line(s) 12833; it names a standard fact of the same chapter that the counted argument uses, and every step of the counted proof built on it is retained verbatim in the delivered body.\par\medskip
\ifdefined\hypertarget\hypertarget{sourcepointer-3}{}\fi
\noindent\textit{Source reference.} Original reference 0654: sources/stacks-sites-cohomology/pinned-original-source.tex, source lines 12432--12432. Prerequisite pointer invoked at proof line(s) 12836; it names a standard fact of the same chapter that the counted argument uses, and every step of the counted proof built on it is retained verbatim in the delivered body.\par\medskip
\begin{lemma}
\label{lemma-perfect}
Let $(\mathcal{C}, \mathcal{O})$ be a ringed site.
Let $E$ be an object of $D(\mathcal{O})$.
The following are equivalent
\begin{enumerate}
\item $E$ is perfect, and
\item $E$ is pseudo-coherent and locally has finite tor dimension.
\end{enumerate}
\end{lemma}

\begin{proof}
Assume (1). Let $U$ be an object of $\mathcal{C}$.
By definition there exists a covering $\{U_i \to U\}$ such that
$E|_{U_i}$ is represented by a strictly perfect complex.
Thus $E$ is pseudo-coherent (i.e., $m$-pseudo-coherent for all $m$) by
Lemma \ref{lemma-pseudo-coherent-independent-representative}.
Moreover, a direct summand of a finite free module is flat, hence
$E|_{U_i}$ has finite Tor dimension by
Lemma \ref{lemma-tor-amplitude}. Thus (2) holds.

\medskip\noindent
Assume (2). Let $U$ be an object of $\mathcal{C}$.
After replacing $U$ by the members of a covering
we may assume there exist integers $a \leq b$ such that $E|_U$
has tor amplitude in $[a, b]$. Since $E|_U$ is $m$-pseudo-coherent
for all $m$ we conclude using Lemma \ref{lemma-perfect-precise}.
\end{proof}

\end{document}
